\begin{prob}{011}{}{2023年11月号 学コン6番}
 実部と虚部がともに整数であるような複素数全体の集合を$A$で表す. また, $a=1+i$とおく(ただし$i=\sqrt{-1}$). 以下の問に答えよ.
 \begin{enumerate}
  \item $A$の元$z_1,z_2$に対し, $\dfrac{z_1z_2}{a}\in A$であるとする. このとき, $\dfrac{z_1}{a}$と$\dfrac{z_2}{a}$のいずれか一方は$A$に属すことを証明せよ.
  \item $x^3 + ay^3 + a^2 z^3 + axyz = 0$を満たす$A$の要素の組$(x,y,z)$を考える.
  \begin{enumerate}
   \item $(x,y,z)=(\alpha,\beta,\gamma)$のとき, $(x,y,z)=(\dfrac{\alpha}{a},\dfrac{\beta}{a},\dfrac{\gamma}{a})$も同式を満たす$A$の要素であることを示せ.
   \item 組$(x,y,z)$をすべて求めよ.
  \end{enumerate}
 \end{enumerate}
\end{prob}

ここに解答を記述。
